% !TEX program = xelatex
\documentclass[12pt,a4paper]{article}


\usepackage{geometry}
\geometry{a4paper,left=2.5cm,right=2.5cm,top=2.5cm,bottom=2.5cm}
\usepackage{ctex}
\usepackage{sectsty}
\sectionfont{\centering}
\renewcommand\thesection{\chinese{section}、}
\renewcommand\thesubsection{\arabic{section}.\arabic{subsection}}
\usepackage[dvipsnames]{xcolor}
\usepackage{booktabs}
\usepackage{makecell}
\usepackage{amsfonts,amsmath,amssymb,mathtools}
\usepackage{multirow}
\usepackage{graphicx}
\usepackage{subcaption}
\usepackage{float}
\usepackage{enumitem}
\usepackage{array}
\usepackage{longtable}
\usepackage{listings}
\usepackage{fancyhdr}
\usepackage{hyperref}
\usepackage{url}

\hypersetup{
  colorlinks=true,
  linkcolor=black,
  urlcolor=blue,
  citecolor=black,
  pdfstartview=Fit,
  breaklinks=true
}

\lstset{
  frame=tb,
  aboveskip=3mm,
  belowskip=3mm,
  showstringspaces=false,
  columns=flexible,
  framerule=1pt,
  rulecolor=\color{gray!35},
  backgroundcolor=\color{gray!5},
  basicstyle={\small\ttfamily},
  numbers=left,
  numberstyle=\tiny\color{gray},
  keywordstyle=\color{blue},
  commentstyle=\color{ForestGreen},
  breaklines=true,
  breakatwhitespace=true,
  tabsize=3
}

\fancyhf{}
\cfoot{\thepage}
\renewcommand{\headrulewidth}{0pt}
\pagestyle{fancy}

\begin{document}

\pagenumbering{arabic}
\setcounter{page}{1}

% ============================================================
% 标题与摘要
% ============================================================
\begin{center}
{\zihao{-2}\bfseries 多火箭残骸音爆定位模型}


\vspace{2ex}
{\zihao{3}\bfseries 摘要}
\end{center}

本文针对多级火箭分离后多个残骸在坠落过程中产生跨音速音爆的定位问题，建立了从“单事件到达时间定位”到“多事件数据关联与联合定位”再到“带随机时间误差的稳健定位”的一整套模型。首先将经纬度按题目给定比例转换为局部平面坐标，把每个音爆事件抽象为位置与时刻均为未知的瞬时声源，得到到达时间方程；再通过时间差（TDOA）消去音爆时刻，将问题化为非线性最小二乘反演。

针对问题1，由“位置 3 个未知量 + 音爆时刻 1 个未知量”得到结论：单个残骸的精确反演至少需要 4 台设备；4 台虽可解，但为排除双解与镜像解、提高稳健性，应使用更多设备。针对 7 台设备的算例，本文先利用两两到达时间差的几何必要条件对数据做一致性筛选，剔除与同一声事件明显不自洽的记录，再采用带多初值的高斯--牛顿法与列文伯格--马夸尔特法求解，得到该残骸音爆位置约为 $110.5000^\circ\mathrm{E}$、$27.3100^\circ\mathrm{N}$，高程约 12514 m，音爆时刻约 12.000 s，7 台设备的到达时间残差均不超过 $3\times10^{-4}$ s。

针对问题2，本文引入每台设备的关联排列 $\sigma_i$，把“哪一个波属于哪一个残骸”建模为离散排列变量，并与 16 个连续参数（4 个残骸的位置与时刻）共同构成混合优化模型。通过自由度分析说明：若波的归属未知，4 台设备只能得到 16 个观测、恰好等于 16 个连续未知量，无法判别错误关联；需要 $4N>16$，即至少 5 台设备，才能利用超定残差唯一确定关联。

针对问题3，本文采用“固定关联—分别定位”与“固定位置—匈牙利算法更新关联”的交替迭代求解，结合“四个音爆时刻相差不超过 5 s”的约束筛选，解出四个残骸的位置、高程与音爆时刻；所给四组解的全部 28 个到达时间残差均小于 $6\times10^{-4}$ s，与表中三位小数的记录精度一致，验证了关联与定位结果的正确性。

针对问题4，将到达时间误差建模为 $U(-0.5,0.5)$ 的随机误差，把问题2的模型修正为加权最小二乘联合优化，并给出蒙特卡洛误差分析：在 7 台设备、关联正确的条件下，500 次试验中四个残骸三维定位误差均值约为 230--730 m；若进一步将台站加密到 12 台并优化布局，总体平均误差约 225 m、95\% 分位数约 532 m，可实现“误差小于 1 km”的目标。

模型的优点在于把“数据关联”与“几何定位”统一在一个目标函数中求解，并用一致性检验、残差分析与蒙特卡洛试验形成完整验证链；通过增加台站数、改善几何布局或引入弹道约束，还可进一步抑制时间误差的影响。

\vspace{2ex}
\noindent{\heiti 关键词：} 音爆定位；最小二乘法非线性规划；匈牙利算法；混合优化；蒙特卡洛误差分析

% ============================================================
% 一、问题重述
% ============================================================
\clearpage
\section{问题重述}

\subsection{问题背景}
多级火箭在飞行过程中会分离下面级火箭和助推器，这些分离的残骸在下落过程中会穿越音速附近，产生跨音速音爆。在理论落区布置监测设备即可记录各残骸音爆到达的时刻。利用各级残骸的到达时间，就能反演残骸发生音爆时的空间位置与时刻。结合弹道外推估计落点即可为快速回收提供依据。

多残骸场景中，单台设备无法得知第几组波来自哪一个残骸，需要结合多台设备来进行分析。

因此，本问题含有两方面的挑战：一是“定位”，即求解残骸音爆的空间位置与时刻；二是“数据关联”，即判别每台设备收到的波属于哪一个残骸。

\subsection{问题重述}

（1）根据物理条件建立数学模型，分析空中单个残骸发生音爆时的经度、纬度、高程与时间所需的最少设备数，并从 7 台设备的数据中选取合适的数据，计算该残骸发生音爆的位置与时间。

（2）假设空中有 4 个残骸，每台设备按时间先后收到 4 组震动波。建立模型确定每组波属于哪一个残骸，并回答确定 4 个残骸位置与时间所需的最少设备数。

（3）利用 7 台设备、每台 4 组到达时间的数据，解出 4 个残骸的位置、高程与音爆时刻，其中四个音爆时刻相差不超过 5 s。

（4）假设设备记录时间存在 0.5 s 的随机误差，修正问题2的模型并给出带误差算例与误差分析；若时间误差无法降低，提出实现误差小于 1 km 的定位方案，并依据问题3结果模拟布站与观测数据加以验证。

% ============================================================
% 二、问题分析
% ============================================================
\section{问题分析}

\subsection{问题1的分析}

单个残骸发生音爆时，位置需要 3 个坐标描述，加上 1 个音爆时刻，共 4 个未知量。每台设备只提供 1 个到达时刻，因此从自由度上至少需要 4 台设备。但到达时间方程是非线性的，4 台设备在数学上可能出现两个镜像解，且测量误差会被放大；因此理论下限是 4 台，实际求解应采用更多设备形成超定方程组。

对于给出的 7 台数据，先利用“同一事件在两台设备上的时间差不可能超过两台设备间距除以声速”这一必要条件做数据体检，再从一致的数据中选择/组合求解，最后用全部台站残差交叉验证，避免个别异常记录污染定位结果。

\subsection{问题2的分析}

问题2在问题1的基础上增加“数据关联”这一离散层：每台设备收到的 4 组波按时间先后排列，但不同残骸到不同设备的传播距离不同，所以“先听到”不代表“先发生”，同一残骸在各设备中的到达顺序也可能不同。不能简单地把所有设备第 $k$ 早的波归为一组。

因此把关联变量建模为排列，将离散关联与连续位置时刻放入同一个残差平方和目标函数。求解时采用“固定关联—定位”与“固定位置—指派”的交替迭代。设备数方面，需要同时比较未知量数与观测数：4 台设备时观测数恰等于连续未知量数，无法判别关联，5 台设备才使系统超定。

\subsection{问题3的分析}

问题3给出 7 台设备、每台 4 个到达时间。应先恢复“哪个波属于哪个残骸”的关联矩阵，再对每个残骸分别执行问题1的定位。利用四个音爆时刻相差不超过 5 s、位置落在合理空域等物理约束可以大幅缩小候选关联的搜索范围，最终以“28 个到达时间全部自洽”作为验证标准。

\subsection{问题4的分析}

0.5 s 的时间误差折算为距离约 170 m，会同时影响定位精度与关联判断。因此需要把问题2模型中的普通最小二乘修正为加权最小二乘，并在关联层使用残差假设检验；通过蒙特卡洛试验统计位置误差。若时间误差无法进一步降低，则应通过增加台站、优化几何布局、引入弹道与四残骸之间的时间约束等方式换取空间冗余，从而实现 1 km 精度。

% ============================================================
% 三、模型假设
% ============================================================
\section{模型假设}

\begin{enumerate}[label=(\arabic*)]
  \item 声速恒定，取 $c=340$ m/s，忽略声速随温度、高度与风速的变化；
  \item 计算两点间距离时忽略地面曲率，经度每度按 $97.304$ km、纬度每度按 $111.263$ km 换算为平面距离；
  \item 每个残骸的一次音爆可视为一个瞬时声源事件，声波沿直线以球面波形式传播；
  \item 各台监测设备时钟已同步，到达时间记录相互可比；
  \item 第 $i$ 台设备收到的 4 组震动波恰好来自 4 个不同的残骸，即每个残骸在该设备上只产生一个到达记录；
  \item 在问题4中，各设备记录时间误差相互独立，误差上限为 0.5 s；
  \item 残骸音爆点在落区附近合理空域内，音爆时刻不晚于其对应的到达时刻。
\end{enumerate}

% ============================================================
% 四、符号说明
% ============================================================
\section{符号说明}

\begin{table}[H]
\centering
\caption{主要符号及其说明}
\begin{tabular}{clp{7.0cm}}
\toprule
符号 & 含义 & 说明 \\
\midrule
$i$ & 监测设备编号 & $i=1,\dots,N$，本题 $N=7$ \\
$d$ & 残骸编号 & $d=1,2,3,4$ \\
$S_i=(x_i,y_i,h_i)$ & 第 $i$ 台设备坐标 & 经纬度已换算为平面坐标，$h_i$ 为高程 \\
$P_d=(x_d,y_d,h_d)$ & 残骸 $d$ 音爆点位置 & 未知量 \\
$\tau_d$ & 残骸 $d$ 的音爆时刻 & 未知量 \\
$t_{i,k}$ & 第 $i$ 台设备第 $k$ 早到达时刻 & 已知量，$k=1,\dots,4$ \\
$t_{i,\sigma_i(d)}$ & 分配给残骸 $d$ 的实测到达时刻 & $\sigma_i(d)$ 为顺序号 \\
$\sigma_i$ & 第 $i$ 台设备的关联排列 & 双射 $\{1,2,3,4\}\to\{1,2,3,4\}$ \\
$y_{i,k,d}$ & 0-1 关联变量 & 第 $k$ 早波来自残骸 $d$ 时取 1 \\
$c$ & 声速 & $c=340$ m/s \\
$\varepsilon$ & 记录时间随机误差 & $\varepsilon\sim U(-0.5,0.5)$ \\
\bottomrule
\end{tabular}
\end{table}

% ============================================================
% 五、模型建立与求解
% ============================================================
\section{模型的建立与求解}
 
\subsection{公共几何框架}

以 $(110^\circ\mathrm{E},\ 27^\circ\mathrm{N})$ 为参考点，把经度、纬度差转换为米制平面坐标：
\[
x_i=(\lambda_i-110)\times 97304,\qquad
y_i=(\varphi_i-27)\times 111263.
\]
若残骸 $d$ 在 $P_d$ 于时刻 $\tau_d$ 发生音爆，则第 $i$ 台设备到达时间为
\[
t_{id}=\tau_d+\frac{\|P_d-S_i\|}{c},
\]
其中
\[
\|P_d-S_i\|=\sqrt{(x_d-x_i)^2+(y_d-y_i)^2+(h_d-h_i)^2}.
\]
等价地，有平方方程
\[
\|P_d-S_i\|^2=c^2\left(t_{id}-\tau_d\right)^2.
\]

将两台设备方程相减可消去 $\tau_d$，得到到达时间差方程：
\[
\|P_d-S_i\|-\|P_d-S_r\|=c\left(t_{id}-t_{rd}\right).
\]
该式不显含音爆时刻，适合求初值与做一致性检验。

\subsection{问题1：单个残骸定位模型}

\subsubsection{最少设备数}

未知量为 $x_d,y_d,h_d,\tau_d$ 共 4 个。每台设备给出 1 条到达时间方程，因此理论上最少需要 4 台设备。若音爆时刻已知，则只需 3 台，但本题时刻未知。

应当说明：4 台只保证方程数与未知数相同，可能出现两个镜像解；为保证唯一性和稳健性，实际建议使用不少于 5 台，或借助物理约束（高程为正、位置在落区附近、音爆时刻小于到达时刻）剔除错误分支。本题提供 7 台设备，可构造超定方程组以最小二乘求解。

\subsubsection{超定定位模型}

取 7 台设备的观测，定义残差
\[
r_i=\|P-S_i\|-c(t_i-\tau),
\]
求解非线性最小二乘问题
\[
\min_{x,y,h,\tau}\ \sum_{i=1}^{7}r_i^2.
\]
求解时先由任意两台设备的到达时间差方程线性化得到初值，再用高斯--牛顿法或列文伯格--马夸尔特法迭代精化，并从多组初值出发取全局最优。

\subsubsection{数据一致性检验与求解结果}

同一音爆事件必须满足：对任意两台设备 $i,j$，
\[
c|t_i-t_j|\le\|S_i-S_j\|.
\]
按题面表1中的原始坐标逐对检验可发现，若干台站组合（如 $D$--$E$、$E$--$F$）不满足该必要条件；而表1各台到达时间与表3中同一事件的到达时间完全相同，且改用表3中的对应台站坐标后，全部 7 站残差均小于 $3\times10^{-4}$ s。因此本文判定表1中部分台站坐标在转录时与表3不一致，数值求解采用与问题3一致、可复现的坐标表；若严格按表1原始坐标使用全部数据，则不存在满足单声事件物理约束的解。

经多初值非线性最小二乘求解，得到一级残骸音爆事件的结果如表 \ref{tab:q1} 所示。

\begin{table}[H]
\centering
\caption{问题1：一级残骸音爆位置与时刻求解结果}
\label{tab:q1}
\begin{tabular}{lcccc}
\toprule
项目 & 经度 $(^{\circ})$ & 纬度 $(^{\circ})$ & 高程 (m) & 音爆时刻 (s) \\
\midrule
求解结果 & 110.5000 & 27.3100 & 12513.9 & 11.9999 \\
\bottomrule
\end{tabular}
\end{table}

将解代回 7 台设备，预测到达时间与实测到达时间的残差均小于 $3\times10^{-4}$ s，与表中三位小数的记录精度一致，说明求解结果可靠。

\subsection{问题2：多残骸数据关联与联合定位模型}

\subsubsection{为什么不能按“第k早”分组}

第 $i$ 台设备记录到的 4 个到达时刻按时间排序为 $t_{i,1}<t_{i,2}<t_{i,3}<t_{i,4}$。但到达顺序由“音爆时刻＋传播时间”共同决定，不同残骸的位置可以相差数十公里，因此同一个残骸在不同设备中的到达顺序可能不同，不能把“第 $k$ 列”直接当作同一个残骸。

\subsubsection{关联排列与 0-1 变量}

定义第 $i$ 台设备的关联排列
\[
\sigma_i:\{1,2,3,4\}\to\{1,2,3,4\},
\]
其中 $\sigma_i(d)$ 表示残骸 $d$ 的波在第 $i$ 台设备的到达顺序号。因为每台设备恰好收到 4 组波、对应 4 个残骸，$\sigma_i$ 必须是双射。

若关联正确，则
\[
\|P_d-S_i\|=c\left(t_{i,\sigma_i(d)}-\tau_d\right).
\]

等价地，可定义 0-1 变量 $y_{i,k,d}$，并施加约束
\[
\sum_{d=1}^{4}y_{i,k,d}=1,\qquad
\sum_{k=1}^{4}y_{i,k,d}=1.
\]

\subsubsection{联合目标函数}

建立混合优化模型
\[
\min_{\{P_d,\tau_d\},\{\sigma_i\}}\ 
\sum_{i=1}^{N}\sum_{d=1}^{4}
\left[
\|P_d-S_i\|-c\left(t_{i,\sigma_i(d)}-\tau_d\right)
\right]^2,
\]
附加约束
\[
h_d>0,\qquad \tau_d<\min_i t_{i,\sigma_i(d)},
\]
若使用问题3信息，还要求
\[
\max_{d}\tau_d-\min_{d}\tau_d\le 5.
\]

\subsubsection{求解算法}

采用两层交替迭代：

\begin{enumerate}[label=\arabic*.]
  \item 固定当前关联，将数据分成 4 组，对每个残骸求解单源非线性最小二乘问题；
  \item 固定位置与时刻，计算每台设备对 4 个残骸的预测到达时刻 $\hat t_{i,d}$；
  \item 对每台设备求解 4 阶指派问题
  \[
  \min_{\sigma_i}\sum_{d=1}^{4}\left(\hat t_{i,d}-t_{i,\sigma_i(d)}\right)^2,
  \]
  用匈牙利算法获得最优排列；
  \item 重复至残差平方和收敛；
  \item 从多组初始关联出发重复求解，保留总残差最小且满足物理约束的解。
\end{enumerate}

搜索过程中可用必要条件剪枝：
\[
c|t_{i,k}-t_{j,l}|\le\|S_i-S_j\|.
\]

\subsubsection{最少设备数的论证}

4 个残骸的连续未知量为 $4\times4=16$ 个；$N$ 台设备提供 $4N$ 个到达时刻。

\begin{itemize}
  \item 关联已知、仅做定位：每个残骸 4 个未知数对应 4 条方程，$N=4$ 即可；
  \item 关联未知、需同时判别波的来源：$N=4$ 时 $4N=16$，观测数恰等于连续未知量数，错误关联也可构造零残差解，无法唯一判别；
  \item $N=5$ 时 $4N=20>16$，系统超定，只有正确关联能使全部方程同时近似成立。
\end{itemize}

因此本题意义下最少设备数为 5 台。

\subsection{问题3：四个残骸的数据关联与定位}

\subsubsection{关联矩阵}

以第 $A$ 台设备的到达顺序作为残骸编号顺序，经交替迭代与匈牙利指派，得到每台设备中四个残骸的到达顺序号如表 \ref{tab:assoc} 所示。表中每一列构成 $\{1,2,3,4\}$ 的一个排列，说明每台设备恰好收到四个残骸各一次。

\begin{table}[H]
\centering
\caption{问题3数据关联矩阵：表中数字为该残骸在对应设备中的到达顺序号}
\label{tab:assoc}
\begin{tabular}{lcccc}
\toprule
设备 & 残骸1 & 残骸2 & 残骸3 & 残骸4 \\
\midrule
A & 1 & 2 & 3 & 4 \\
B & 2 & 3 & 1 & 4 \\
C & 4 & 3 & 1 & 2 \\
D & 4 & 2 & 3 & 1 \\
E & 3 & 2 & 1 & 4 \\
F & 4 & 2 & 3 & 1 \\
G & 2 & 1 & 4 & 3 \\
\bottomrule
\end{tabular}
\end{table}

例如设备 $D$ 中“残骸4”排第 1 个到达，而残骸1排第 4 个到达，说明四个残骸在各台设备的到达顺序并不相同，直接按列分组是不正确的。

\subsubsection{定位结果}

按表 \ref{tab:assoc} 分组后，对每个残骸用 7 台设备做最小二乘定位，结果如表 \ref{tab:q3} 所示。

\begin{table}[H]
\centering
\caption{问题3：四个残骸音爆位置与时刻}
\label{tab:q3}
\begin{tabular}{lcccc}
\toprule
残骸 & 经度 $(^{\circ})$ & 纬度 $(^{\circ})$ & 高程 (m) & 音爆时刻 (s) \\
\midrule
残骸1 & 110.5000 & 27.3100 & 12513.9 & 11.9999 \\
残骸2 & 110.3000 & 27.6500 & 11477.9 & 14.0000 \\
残骸3 & 110.7000 & 27.6500 & 13468.2 & 15.0000 \\
残骸4 & 110.5000 & 27.9500 & 11528.9 & 13.0014 \\
\bottomrule
\end{tabular}
\end{table}

四个音爆时刻的最小值与最大值分别约为 12.00 s 与 15.00 s，极差约 3.00 s，满足题设“相差不超过 5 s”的约束。

\subsubsection{残差验证}

将四个解代回全部 28 个到达时刻，各残骸的最大时间残差如表 \ref{tab:q3res} 所示。

\begin{table}[H]
\centering
\caption{问题3：各残骸的 7 台最大到达时间残差}
\label{tab:q3res}
\begin{tabular}{lcccc}
\toprule
残骸 & 残骸1 & 残骸2 & 残骸3 & 残骸4 \\
\midrule
最大残差 $(10^{-4}\ \mathrm{s})$ & 2.8 & 5.1 & 2.3 & 2.1 \\
\bottomrule
\end{tabular}
\end{table}

残差均小于 $6\times10^{-4}$ s，对应距离误差小于 0.2 m，说明关联矩阵与定位结果同时正确。

\subsection{问题4：带随机时间误差的稳健定位模型}

\subsubsection{误差模型}

设观测到达时间为真实到达时间与随机误差之和：
\[
\tilde{t}_{i,k}=t_{i,k}+\varepsilon_{i,k}.
\]
本文将“0.5 s 的随机误差”建模为
\[
\varepsilon_{i,k}\sim U(-0.5,\ 0.5),
\]
即每次记录的最大误差不超过 0.5 s，等价的距离噪声约 170 m。

\subsubsection{加权最小二乘修正模型}

定义预测到达时间
\[
\hat{t}_{i,d}=\tau_d+\frac{\|P_d-S_i\|}{c}.
\]
第 $i$ 台设备中分配给残骸 $d$ 的实测值为 $\tilde{t}_{i,\sigma_i(d)}$，修正后的联合目标函数为
\[
\min_{\{P_d,\tau_d\},\{\sigma_i\}}\ 
\sum_{i=1}^{N}\sum_{d=1}^{4}
\left[
\tau_d+\frac{\|P_d-S_i\|}{c}-\tilde{t}_{i,\sigma_i(d)}
\right]^2.
\]

若各台设备精度不同，可取权重
\[
w_{i,d}=\frac{1}{\sigma_{i,d}^{2}},
\]
将上式改为加权残差平方和。由于误差均匀分布时最大似然与平方损失并不严格等价，本文按以下顺序处理：先以普通最小二乘求初值并估计残差尺度，再以加权最小二乘精化；若某台数据残差明显大于 0.5 s 对应量级，则用 Huber 权重或 RANSAC 策略削弱野值影响。

\subsubsection{关联层的稳健处理}

0.5 s 误差使“零残差”不再成立，关联层应采用以下判据：
\begin{enumerate}[label=\arabic*.]
  \item 对每台设备用预测到达时刻与实测到达时刻构造代价矩阵，用匈牙利算法更新排列；
  \item 比较候选关联的加权残差平方和，利用残差量级判断哪一关联与噪声水平相符；
  \item 利用四个音爆时刻相差不超过 5 s 的约束排除不合理关联；
  \item 对可疑候选做留一站交叉验证，剔除使某台设备残差系统增大的错误关联。
\end{enumerate}

\subsubsection{蒙特卡洛误差分析}

以问题3求解结果作为“真值”，对每条到达时间叠加 $U(-0.5,0.5)$ 的误差，重复 500 次试验并统计三维定位误差。为突出时间误差对几何反演的影响，该组统计在完成关联后给出；关联层的正确性单独以残差判据检验。

\begin{table}[H]
\centering
\caption{问题4蒙特卡洛结果（7 台设备，500 次，关联给定）}
\label{tab:q4mc}
\begin{tabular}{lccccc}
\toprule
指标 & 残骸1 & 残骸2 & 残骸3 & 残骸4 & 总体 \\
\midrule
三维误差均值 (m) & 588 & 395 & 230 & 726 & 485 \\
三维误差 95\% 分位 (m) & 1328 & 848 & 469 & 1621 & 1276 \\
音爆时刻误差均值 (s) & 0.63 & 0.32 & 0.28 & 0.82 & --- \\
\bottomrule
\end{tabular}
\end{table}

可见：在 7 台设备与 ±0.5 s 误差下，单次定位三维误差的均值约为 0.2--0.7 km，但 95\% 分位可能超过 1 km，说明原 7 台布局对“误差恒小于 1 km”的目标余量不足。

\subsubsection{时间误差不可降低时的布站优化方案}

若时间误差无法进一步降低，本文采用“增加台站数 + 优化几何布局 + 利用物理约束”的组合方案：

\begin{enumerate}[label=\arabic*.]
  \item 将台站由 7 台增加到 12 台，覆盖经度 $110.05^\circ$--$110.90^\circ$、纬度 $27.10^\circ$--$28.00^\circ$ 的区域，使目标位于台站阵列内部而非一侧；
  \item 台站高程取 500--1000 m 的起伏布置，改善高程方向的观测几何；
  \item 利用四个残骸音爆时刻相差不超过 5 s 的耦合约束进行联合估计；
  \item 可选地引入弹道约束，把各残骸音爆点约束到其下落轨迹附近，进一步抑制随机误差。
\end{enumerate}

依据问题3求得的四个定位结果模拟 12 台设备的到达时间并叠加 ±0.5 s 误差，再做 500 次蒙特卡洛试验，结果如表 \ref{tab:q4new} 所示。

\begin{table}[H]
\centering
\caption{12 台优化布站后的三维定位误差（500 次蒙特卡洛）}
\label{tab:q4new}
\begin{tabular}{lccccc}
\toprule
指标 & 残骸1 & 残骸2 & 残骸3 & 残骸4 & 总体 \\
\midrule
三维误差均值 (m) & 288 & 178 & 182 & 249 & 225 \\
三维误差 95\% 分位 (m) & 705 & 366 & 380 & 574 & 532 \\
\bottomrule
\end{tabular}
\end{table}

试验结果表明，优化布站后总体平均误差约 225 m、95\% 分位数约 532 m，稳定满足“误差小于 1 km”的目标，验证了所提布站与定位方案的可行性。

% ============================================================
% 六、模型检验与误差分析
% ============================================================
\section{模型检验、误差分析与稳健性讨论}

\subsection{残差检验}

问题1与问题3的求解均以“所有到达时间残差接近记录精度”为验收标准。问题3中 28 个残差全部小于 $6\times10^{-4}$ s，且问题1的单个事件亦满足该标准，说明模型方程、坐标转换与关联分组自洽。

\subsection{误差传播分析}

单条到达时间误差 0.5 s 对应
\[
c\times0.5=170\ \mathrm{m}
\]
的距离误差。双曲线定位中，最终定位误差与设备几何有关，可用几何精度因子（GDOP）刻画：台站包围目标、孔径越大，几何放大越小；台站近似共线或目标在阵列外侧时，误差会被显著放大。这正是 12 台优化布站方案能显著压缩误差的原因。

\subsection{灵敏度与稳健性}

对问题4模型，可从以下角度做稳健性检验：
\begin{enumerate}[label=\arabic*.]
  \item 改变误差幅度（如 0.2、0.5、0.8 s），观察定位误差近似线性增长；
  \item 改变随机种子重复试验，检验统计量是否稳定；
  \item 删除任意一台设备后重新求解，检验残差与结果是否保持稳定；
  \item 对关联层的初始分组加扰动，检验交替迭代是否收敛到同一关联。
\end{enumerate}

% ============================================================
% 七、模型评价与推广
% ============================================================
\section{模型的评价与推广}

\subsection{模型的优点}

\begin{enumerate}[label=\arabic*.]
  \item 将“数据关联”与“几何定位”统一为混合优化问题，避免了把不同设备第 $k$ 早的波强行归为一类的错误；
  \item 采用两层交替迭代与匈牙利算法，把离散排列搜索与连续非线性反演分离，求解稳定、可解释；
  \item 利用几何必要条件和物理约束进行数据筛选与剪枝，能够识别异常台站或转录错误；
  \item 通过蒙特卡洛试验给出定位误差分布，并验证了增加台站与优化布站对精度的提升效果。
\end{enumerate}

\subsection{模型的缺点}

\begin{enumerate}[label=\arabic*.]
  \item 模型将声速取为常数，未考虑声速随高度变化、风场和温度梯度引起的传播路径弯曲；
  \item 交替迭代可能收敛到局部最优，理论上的全局收敛性依赖多初值与剪枝策略；
  \item 问题4中的蒙特卡洛误差统计在关联给定条件下进行，实际工程中关联错误会进一步放大误差。
\end{enumerate}

\subsection{模型推广}

该框架可推广到地震震源定位、声音/闪电多站时差定位、通信到达时间定位等具有“未知时刻的多点观测”特征的场景；引入波形特征、到达方向或运动学约束后，可进一步提高关联可靠性与定位精度。

% ============================================================
% 参考文献
% ============================================================
\section{参考文献}

\begin{enumerate}[label={[\arabic*]}]
  \item 2024年深圳杯数学建模挑战赛 A 题：多个火箭残骸的准确定位.
  \item 王正明，易东云. 测量数据建模与参数估计[M]. 长沙：国防科技大学出版社，1996.
  \item 边肇祺，张学工. 模式识别[M]. 北京：清华大学出版社，2000.
  \item Kuhn H W. The Hungarian method for the assignment problem[J]. Naval Research Logistics Quarterly, 1955, 2(1--2): 83--97.
  \item Chan Y T, Ho K C. A simple and efficient estimator for hyperbolic location[J]. IEEE Transactions on Signal Processing, 1994, 42(8): 1905--1915.
\end{enumerate}

% ============================================================
% AI 工具使用声明（如比赛要求请保留；自行按要求修改）
% ============================================================
\section*{AI 工具使用声明}
本文在资料检索、语言组织与代码调试过程中使用了 AI 辅助工具。所有数学模型、公式推导、数值结果与结论均由作者核验，AI 未代替作者完成建模决策与最终修改。

% ============================================================
% 附录
% ============================================================
\appendix
\section{支撑材料说明}

支撑材料包括：
\begin{itemize}
  \item 主程序与绘图脚本（Python）；
  \item 问题1--问题4的数值输出；
  \item 蒙特卡洛误差分析脚本与结果表。
\end{itemize}

\section{题目数据表}

\begin{table}[H]
\centering
\caption{问题3各设备到达时间（单位：s）}
\label{tab:data3}
\begin{tabular}{lcccc}
\toprule
设备 & 第1组 & 第2组 & 第3组 & 第4组 \\
\midrule
A & 100.767 & 164.229 & 214.850 & 270.065 \\
B & 92.453 & 112.220 & 169.362 & 196.583 \\
C & 75.560 & 110.696 & 156.936 & 188.020 \\
D & 94.653 & 141.409 & 196.517 & 258.985 \\
E & 78.600 & 86.216 & 118.443 & 126.669 \\
F & 67.274 & 166.270 & 175.482 & 266.871 \\
G & 103.738 & 163.024 & 206.789 & 210.306 \\
\bottomrule
\end{tabular}
\end{table}

表中每行数值已按时间从小到大排列。本文求解得到的关联矩阵见表 \ref{tab:assoc}。

\section{建议在终稿中补入的图表}

\begin{itemize}
  \item 监测台站与四个残骸位置的三维散点图；
  \item 各残骸音爆到达时刻残差柱状图；
  \item 12 台优化布站俯视图与 GDOP 分布图；
  \item 问题4误差幅度—定位误差灵敏度曲线。
\end{itemize}

\end{document}

\end{document}
